\subsection{Explicit loops, Fermat covers, and the Ihara series}
\paragraph{Statement A9 (Proved; Part A(f)).}\label{res:A9}
For $U=\mathbb P^1_{\mathbb Q}\setminus\{0,1,\infty\}$ choose the rational tangential base point at $0$ directed toward $1$, the corresponding section $s$, and geometric generators $x,y$ around $0,1$. Then
\[
 C_\pi(s(\Gamma_{\mathbb Q}^{\circ}))=\overline{\langle x\rangle},
 \qquad y\notin C_\pi(s(\Gamma_{\mathbb Q}^{\circ})),
 \qquad [x,y]\notin C_\pi(s(\Gamma_{\mathbb Q}^{\circ})).
\]
These assertions concern actual elements and the genuine action defined by $s$. They test the formula in cote 149, section 3, manuscript p.~27, scans 55--56, and the characterizations in Longue Marche, section 34, scan 483.

\begin{proof}
The equality is A6. The image of $y$ in $\pi^{\rm ab}=\widehat{\mathbb Z}(1)^2$ is $(0,1)$, whereas the image of $\overline{\langle x\rangle}$ is the first summand. This proves the assertion for $y$ with the quotient map specified. We give an independent arithmetic obstruction for the commutator.

Let $H_n$ be the kernel of $\pi\to(\mathbb Z/n)^2$ sending $x,y$ to the coordinate generators. Its geometric covering is
\[
 w^n+v^n=1,\qquad z=w^n.
\]
It is finite etale over $U$, and its compactification is the Fermat curve $X_n$. The maps
\[
 \overline{[\pi,\pi]}\longrightarrow H_n\longrightarrow H_n^{\rm ab}
 \longrightarrow T(J_{X_n})
 \tag{A.3}
\]
are continuous and are equivariant for the restricted section over $\mathbb Q(\mu_\infty)$. A finite field extension supplies any chosen compatible lift and does not affect the invariant-vanishing input A1. The first arrow is an inclusion, the second is abelianization, and the last is induced by the open immersion into the compactification. In particular the last two arrows kill the second derived subgroup of $\pi$.

Here is a direct certificate that $[x,y]$ has nonzero image for $n=3$. The abelian covering of the topological wedge of the two generators has vertices $(a,b)\in(\mathbb Z/3)^2$ and directed edges $e_x(a,b),e_y(a,b)$ incrementing one coordinate. Its commutator cycle is
\[
 e_x(0,0)+e_y(1,0)-e_x(0,1)-e_y(0,0).
\]
Compactification kills the cycles around $x^3$, $y^3$ and $(xy)^3$, at all their lifts. Fix a primitive cube root $\omega$ and give $e_x(a,b)$ weight $\omega^{a+b}$ and $e_y(a,b)$ weight zero. This functional vanishes on all the boundary cycles: the relevant sums are $1+\omega+\omega^2$ and $1+\omega^2+\omega^4$. Its value on the commutator is $1-\omega\ne0$. The calculation works over $\mathbb Z_3[\omega]$, so it proves nonvanishing in $T_3(J_{X_3})$, through the comparison of topological finite covers with etale finite covers. The exact source for that comparison is SGA~1, Expose XII, Theorem 5.1: finite etale covers of a complex scheme locally of finite type correspond to its finite analytic etale covers \cite{SGA1}. The map on first homology is precisely the pro-$3$ abelianization of the displayed cover.

An element fixed by $s(\Gamma_{\mathbb Q}^{\circ})$ would have fixed image in this Tate module, contrary to A1. This proves the commutator assertion without using its conclusion as an input to the Fermat calculation.

To identify the requested Ihara power series, use the free pro-$\ell$ quotient $P$ of $\pi$, $M=P'/P''$, and $\Lambda=\mathbb Z_\ell[[T_1,T_2]]$. The opened primary research source is Kodani's thesis, section 4.3, Lemma 4.3.1 and equation (4.3.2), which restates Ihara's Theorem 2(i) and proves the module assertion independently \cite{Kodani2017,Ihara1991}. Ihara's own opened ICM paper, section 6.4, formula (6.4.1), gives the corresponding multiplier $B'_\sigma$ and the Fermat-Jacobian connection on pp.~116--117. In the thesis's tangential convention, $M$ is free of rank one over $\Lambda$, generated by the commutator, and
\[
 \rho(g)([a,b])=F_g(T_1,T_2)[a,b]\quad\text{in }M.
 \tag{A.4}
\]
The action on $\Lambda$ sends $1+T_i$ to $(1+T_i)^{\chi_\ell(g)}$. The paragraph following Theorem 4.3.3 identifies the resulting module with the inverse limit of the Tate modules of the Fermat Jacobians; its transition maps are induced by the power maps between Fermat curves. Equation (A.3) is the finite-level compactification map in that construction. The exact hypotheses here are $\{0,1,\infty\}\subset\mathbb P^1_{\mathbb Q}$, a fixed prime $\ell$, the pro-$\ell$ geometric group, and the specified tangential Galois action.

We also check the change of tangential convention. Applying $z\mapsto1/z$ to the convention with tangent at infinity takes it to the present tangent at zero. In the old generators $a,b$, the new generators are $x=(ab)^{-1}$ and $y=b$. The identity
\[
 [x,y]=x[a,b]^{-1}x^{-1}
\]
shows that their commutator bases differ by the invertible scalar $-(ab)^{-1}$ in the group ring. On $\Gamma_{\mathbb Q}^{\circ}$ the scalar ring has trivial action, so the multiplier $F_g$ in (A.4) is unchanged. Taking $\ell=3$, the preceding nonzero Fermat image and A1 show that $F_g\ne1$ for at least one $g\in\Gamma_{\mathbb Q}^{\circ}$. Otherwise (A.4) would fix the commutator and hence its nonzero image under (A.3). This supplies the requested power-series obstruction with its actual module and comparison maps; no numerically guessed Galois matrix is used.
\end{proof}

The executed program \texttt{code/root\_math\_checks.py} computes the incidence and boundary matrices for $n=2,3,4,5$. Their compact first-homology ranks are $0,2,6,12$; the commutator raises the boundary rank exactly for $n=3,4,5$. The degree-three calculation above is the integral proof behind these checks. The CM elliptic test, including the image of the full centralizer, is given in E1--E3 below.
